%%═══════════════════════════════════════════════════════════════════
%% Lecture 02 — Classical Field Theory (continued):
%%              Euler–Lagrange for Fields, Examples,
%%              Relativistic Notation, Natural Units
%% 77609 — Introduction to Elementary Particles
%% Date: 26/03/2026 
%% Lecturer: Prof. Yonit Hochberg
%%═══════════════════════════════════════════════════════════════════

\section{Lecture 02 — Euler--Lagrange Equations for Fields, Basic
  Examples, and Relativistic Notation}\label{sec:lec02}
\begin{center}
\begin{tabular}{ll}
  \textbf{Date:}\quad 26/03/2026   & \\
\end{tabular}
\end{center}
\hrule\vspace{1.5em}
%%──────────────────────────────────────────────────────────────────
\subsection*{Overview}

This lecture continues the construction of \emph{classical field theory}
as the direct generalization of analytical mechanics in which the
\textbf{generalized coordinates are fields} \(\phi(x,t)\) rather than
particle positions \(x(t)\). Starting from an \emph{action}
\(S=\int \lagr\,dx\,dt\) built from a \emph{local} (single spacetime point)
and \emph{real} Lagrangian density \(\lagr\), we impose the principle of
minimal action (\(\delta S=0\)) to obtain the Euler--Lagrange equations
for fields. Several instructive examples are worked out: the wave
equation, the Klein--Gordon equation (adding a ``mass term''), and a
nonlinear interacting equation from a cubic interaction \(\sim\phi^3\),
which motivates perturbation theory.

The second half of the lecture is a compact but thorough recap of
\textbf{relativistic notation} needed for Lorentz-invariant field
theories: four-vectors, upper/lower indices, Einstein summation, the
mostly-negative metric \(g_{\mu\nu}\), raising and lowering indices, the
Lorentz-invariant scalar product, four-derivatives \(\partial_\mu\), and
the d'Alembertian (box) operator \(\Box\). These tools are then used to
rewrite the earlier Lagrangians and equations of motion in manifestly
Lorentz-covariant form. The lecture closes with a brief discussion of
the Lorentz group condition \(\Lambda^T g \Lambda = g\), the
classification timelike/lightlike/spacelike via \(A^2\), the
transformation of tensors, and the convention of \textbf{natural units}
\(\hbar=c=1\).

%%──────────────────────────────────────────────────────────────────
\subsection*{Foundational viewpoint: fields as generalized coordinates}

\begin{remark}[The conceptual shift]
In classical mechanics we solve for a finite number of generalized
coordinates \(q(t)\) (e.g.\ \(x(t)\)). In classical field theory the
generalized coordinate is a \emph{field} \(\phi(x,t)\); solving the theory
means determining \(\phi\) for \emph{all} spacetime points \((x,t)\). In
\(d\) spatial dimensions, \(\phi\) depends on \(d+1\) arguments.
\end{remark}

%%──────────────────────────────────────────────────────────────────
\subsection*{Axioms/assumptions for the classical field theories we study}

We consider a theory with (possibly) several fields \(\phi_i(x,t)\), but
for most of the lecture we restrict to a \textbf{single real scalar
field} \(\phi(x,t)\) in \textbf{\(1+1\) dimensions} (one space coordinate
\(x\) and time \(t\)).

\dfn{Action and Lagrangian density}{The action functional is
\begin{equation}
  S[\phi] = \int dt\,dx\; \lagr\big(\phi, \partial_t\phi, \partial_x\phi; x,t\big),
\end{equation}
where \(\lagr\) is the \emph{Lagrangian density} (denoted \(\mathcal{L}\)
in lecture).}

The lecture adopts simplifying assumptions that mirror standard choices
in mechanics:
\begin{itemize}
  \item \textbf{No explicit spacetime dependence:} we focus on
    \(\lagr=\lagr(\phi,\dot\phi,\phi')\) with no explicit \(x\) or \(t\)
    dependence. This is the field-theory analogue of ``no explicit time
    dependence'' in mechanics and encodes \textbf{translational
    invariance}, ensuring conservation of energy and momentum (via
    Noether's theorem).
  \item \textbf{First derivatives only:} we take \(\lagr\) to depend on
    \(\phi\) and at most \(\dot\phi\equiv\partial_t\phi\) and
    \(\phi'\equiv\partial_x\phi\). Higher-derivative theories exist but
    are beyond the baseline framework introduced here.
  \item \textbf{Locality:} \(\lagr\) depends on fields evaluated at the
    \emph{same} spacetime point \((x,t)\). This defines a \textbf{local
    field theory}.
  \item \textbf{Reality:} \(\lagr\in\mathbb{R}\). Physically, this
    ensures a real Hamiltonian and (in the quantum theory) unitary time
    evolution / probability conservation.
\end{itemize}

%%──────────────────────────────────────────────────────────────────
\subsection*{Euler--Lagrange equation for fields}

\thm{Euler--Lagrange equation in \(1+1\) dimensions}{For an action
\begin{equation}
  S[\phi]=\int dt\,dx\;\lagr(\phi,\dot\phi,\phi')
\end{equation}
the principle of minimal action (\(\delta S=0\)) gives the field
equation
\begin{equation}
  \frac{\partial \lagr}{\partial \phi}
  = \frac{\partial}{\partial t}\left(\frac{\partial \lagr}{\partial \dot\phi}\right)
  + \frac{\partial}{\partial x}\left(\frac{\partial \lagr}{\partial \phi'}\right).
  \label{eq:lec02:EL-1p1}
\end{equation}}

\begin{remark}[Analogy with mechanics]
Compare \eqref{eq:lec02:EL-1p1} with
\(\frac{d}{dt}\frac{\partial L}{\partial \dot q}-\frac{\partial L}{\partial q}=0\).
Field theory introduces an additional ``evolution-like'' variable \(x\),
so stationarity of the action produces \emph{both} a \(t\)-derivative term
and an \(x\)-derivative term. This is a concrete sense in which
\textbf{space and time appear on similar footing} in the variational
principle.
\end{remark}

%%──────────────────────────────────────────────────────────────────
\subsection*{Worked examples: free waves, mass term, and nonlinearity}

\subsubsection*{Example 1: the wave equation}

Take
\begin{equation}
  \lagr = \frac12 \dot\phi^{\,2} - v^2(\phi')^2.
  \label{eq:lec02:L-wave}
\end{equation}
Here \(v\) is a constant with units of velocity (the wave speed).

\begin{itemize}
  \item Since \(\lagr\) has no explicit \(\phi\) dependence,
    \(\partial\lagr/\partial\phi = 0\).
  \item \(\partial\lagr/\partial\dot\phi = \dot\phi\), so
    \(\partial_t(\partial\lagr/\partial\dot\phi)=\ddot\phi\).
  \item \(\partial\lagr/\partial\phi'\propto -v^2\phi'\), so the spatial
    contribution yields a term proportional to \(v^2\phi''\).
\end{itemize}

The resulting equation of motion is the \textbf{wave equation}
\begin{equation}
  \ddot\phi = v^2 \phi''.
  \label{eq:lec02:wave-eom}
\end{equation}
A plane-wave solution is
\begin{equation}
  \phi(x,t) = e^{i(kx-\omega t)},
  \qquad
  \omega = kv.
  \label{eq:lec02:wave-sol}
\end{equation}

\begin{remark}[Why this is ``free'']
The equation is linear, so superpositions of solutions are solutions.
There is no term nonlinear in \(\phi\); physically, different waves do
not interact with each other.
\end{remark}

\subsubsection*{Example 2: adding a ``mass term'' \(-m^2\phi^2\)}

Now take
\begin{equation}
  \lagr = \frac12 \dot\phi^{\,2} - v^2(\phi')^2 - m^2 \phi^2,
  \label{eq:lec02:L-KG}
\end{equation}
where the additional term \(-m^2\phi^2\) is called the \textbf{mass term}.

Now \(\partial\lagr/\partial\phi\neq 0\), and the Euler--Lagrange equation
becomes
\begin{equation}
  \ddot\phi = v^2\phi'' - m^2 \phi.
  \label{eq:lec02:KG-eom-1p1}
\end{equation}
This is the \textbf{Klein--Gordon equation} (in the \(1+1\) notation used
in class; later it will be written in a Lorentz-covariant form).

Try the same plane-wave ansatz \(\phi=e^{i(kx-\omega t)}\). Plugging into
\eqref{eq:lec02:KG-eom-1p1} yields the modified dispersion relation
\begin{equation}
  \omega^2 = (kv)^2 + m^2.
  \label{eq:lec02:KG-dispersion}
\end{equation}

\begin{remark}[Physical intuition for the mass term]
Compared to the pure wave equation, the additional \(-m^2\phi\) term acts
like a restoring force even at \(k=0\): the field can oscillate locally
with frequency \(\omega=m\). In the quantum theory, this same parameter
will set the rest energy of the corresponding particle.
\end{remark}

\subsubsection*{Example 3: an interaction \(\sim \phi^3\) and nonlinearity}

Consider
\begin{equation}
  \lagr = \frac12 \dot\phi^{\,2} - v^2(\phi')^2 + \frac{\lambda}{3}\phi^3.
  \label{eq:lec02:L-cubic}
\end{equation}
The Euler--Lagrange equation becomes
\begin{equation}
  \ddot\phi = v^2\phi'' + \lambda \phi^2.
  \label{eq:lec02:cubic-eom}
\end{equation}
This is \textbf{nonlinear} because of the \(\phi^2\) term. Unlike the
previous two examples, there is generally \textbf{no closed-form exact
solution} for arbitrary initial data.

\begin{remark}[Why this matters for the course: perturbation theory]
Interactions corresponding to terms \(\phi^n\) with \(n\ge 3\) typically
produce nonlinear equations. The course strategy will be
\textbf{perturbation theory}: expand around a minimum of an effective
potential and treat higher powers of the field as small corrections.
This will be the seed of the perturbative expansion that later becomes
Feynman diagrams in the quantum theory.
\end{remark}

%%──────────────────────────────────────────────────────────────────
\subsection*{Relativistic notation: four-vectors and indices}

Up to this point we wrote \(x\) and \(t\) explicitly. For most of the
systems of interest in particle physics we want \textbf{Lorentz
invariance} to be manifest, so it is efficient to adopt relativistic
notation.

\subsubsection*{Spacetime interval}

Special relativity implies invariance of
\begin{equation}
  s^2 = c^2t^2 - x_i x_i,
  \qquad i=1,2,3,
\end{equation}
between inertial frames.

\subsubsection*{Four-vectors \(x^\mu\) and general \(A^\mu\)}

\dfn{Position four-vector}{The spacetime coordinate of an event is
\begin{equation}
  x^\mu = (ct,\,x,\,y,\,z),
  \qquad \mu=0,1,2,3.
\end{equation}
All four components have dimensions of length.}

For a boost in the \(x\)-direction,
\begin{align}
  x^{0'} &= \gamma(x^0-\beta x^1), &
  x^{1'} &= \gamma(x^1-\beta x^0), &
  x^{2'} &= x^2, &
  x^{3'} &= x^3,
\end{align}
where \(\beta=v/c\) and \(\gamma = 1/\sqrt{1-\beta^2}\).

\dfn{Four-vector}{A four-vector is any object \(A^\mu=(A^0,A^1,A^2,A^3)\) that transforms
under boosts/rotations in the same way as \(x^\mu\).}

\ex{Four-momentum}{\(p^\mu=(E/c,\bvec{p})\) is a four-vector.}

\subsubsection*{Scalar (dot) product}

\dfn{Lorentz-invariant scalar product}{The scalar product of four-vectors \(A^\mu\) and \(B^\mu\) is defined as
\begin{equation}
  A\cdot B \equiv A^0B^0 - A^1B^1 - A^2B^2 - A^3B^3.
  \label{eq:lec02:dot}
\end{equation}}
With this definition, \(A\cdot B\) is Lorentz invariant.

\subsubsection*{Lower indices and Einstein summation}

\dfn{Lower-index components}{Define the lower-index version of the position four-vector by
\begin{equation}
  x_\mu = (ct,\,-x,\,-y,\,-z),
  \label{eq:lec02:x-lower}
\end{equation}
and similarly for any four-vector \(A^\mu\),
\begin{equation}
  A_\mu = (A^0,\,-A^1,\,-A^2,\,-A^3).
  \label{eq:lec02:A-lower}
\end{equation}}

\dfn{Einstein summation convention}{Any index appearing twice is summed from \(0\) to \(3\).}

Then
\begin{equation}
  A^\mu B_\mu
  = A^0B_0 + A^1B_1 + A^2B_2 + A^3B_3
  = A^0B^0 - A^1B^1 - A^2B^2 - A^3B^3
  = A\cdot B.
  \label{eq:lec02:einstein-dot}
\end{equation}

\begin{remark}[Upper vs.\ lower indices]
Quantities like \(A^\mu B^\mu\) (both indices upstairs) are not Lorentz
invariant and should not appear in physical expressions. In sensible
Lorentz-covariant formulas, repeated indices appear once upstairs and
once downstairs.
\end{remark}

\begin{remark}[Index sanity checks from lecture]
\begin{itemize}
  \item Never repeat an index more than twice in the same product
    (\(A_\mu B^\mu C^\mu\) is ambiguous).
  \item The name of a dummy summed index does not matter:
    \(A_\mu B^\mu = A_\nu B^\nu\).
  \item All indices contracted \(\Rightarrow\) scalar; one free index
    \(\Rightarrow\) vector; two free indices \(\Rightarrow\) rank-2
    tensor; etc.
\end{itemize}
\end{remark}

\subsubsection*{Metric tensor and raising/lowering}

\dfn{Metric tensor}{The Minkowski metric (``mostly negative'') is
\begin{equation}
  g_{\mu\nu}=\mathrm{diag}(1,-1,-1,-1),
  \label{eq:lec02:metric}
\end{equation}
often denoted \(\eta_{\mu\nu}\).}

It lowers indices:
\begin{equation}
  A_\mu = g_{\mu\nu}A^\nu,
  \label{eq:lec02:lowering}
\end{equation}
and raises them with \(g^{\mu\nu}\), which has the same numerical
entries:
\begin{equation}
  A^\mu = g^{\mu\nu}A_\nu,
  \qquad g^{\mu\nu}=\mathrm{diag}(1,-1,-1,-1).
  \label{eq:lec02:raising}
\end{equation}

The scalar product can be written as
\begin{equation}
  A\cdot B = A^\mu B_\mu = A^\mu g_{\mu\nu} B^\nu.
  \label{eq:lec02:dot-metric}
\end{equation}

%%──────────────────────────────────────────────────────────────────
\subsection*{Four-derivatives and the box operator}

\dfn{Four-derivative}{Define
\begin{equation}
  \partial_\mu \equiv \frac{\partial}{\partial x^\mu}
  = \left(\frac{1}{c}\frac{\partial}{\partial t},\,\frac{\partial}{\partial x},\,
  \frac{\partial}{\partial y},\,\frac{\partial}{\partial z}\right),
  \label{eq:lec02:partial-lower}
\end{equation}
and
\begin{equation}
  \partial^\mu = g^{\mu\nu}\partial_\nu
  = \left(\frac{1}{c}\frac{\partial}{\partial t},\,-\frac{\partial}{\partial x},\,
  -\frac{\partial}{\partial y},\,-\frac{\partial}{\partial z}\right).
  \label{eq:lec02:partial-upper}
\end{equation}}

\ex{Derivative of a scalar is a four-vector}{If \(f(x)\) is a scalar function of spacetime, then \(\partial_\mu f\) is
a four-vector.}

\ex{Divergence of a four-vector is a scalar}{If \(h^\mu(x)\) is a four-vector field, then \(\partial_\mu h^\mu\) is a
Lorentz scalar.}

\dfn{D'Alembertian / box operator}{\begin{equation}
  \Box \equiv \partial^\mu\partial_\mu
  = \frac{1}{c^2}\frac{\partial^2}{\partial t^2}
  - \frac{\partial^2}{\partial x^2}
  - \frac{\partial^2}{\partial y^2}
  - \frac{\partial^2}{\partial z^2}.
  \label{eq:lec02:box}
\end{equation}}

%%──────────────────────────────────────────────────────────────────
\subsection*{Rewriting the scalar field theory in covariant form}

The free scalar Lagrangian with a mass term can be written as
\begin{equation}
  \lagr
  = \frac12\,\partial_\mu\phi\,\partial^\mu\phi - m^2\phi^2.
  \label{eq:lec02:L-covariant}
\end{equation}
The \(\partial_\mu\phi\,\partial^\mu\phi\) term is called the
\textbf{kinetic term}, and \(-m^2\phi^2\) is the \textbf{mass term}.

The Klein--Gordon equation becomes
\begin{equation}
  (\Box + m^2)\phi = 0.
  \label{eq:lec02:KG-box}
\end{equation}

The Euler--Lagrange equation in relativistic notation was given as
\begin{equation}
  \partial_\mu\left(\frac{\partial \lagr}{\partial(\partial_\mu\phi)}\right)
  = \frac{\partial \lagr}{\partial\phi}.
  \label{eq:lec02:EL-covariant}
\end{equation}

%%──────────────────────────────────────────────────────────────────
\subsection*{Lorentz group, invariants, and tensors (recap)}

Write a Lorentz transformation as
\begin{equation}
  x^{\mu'} = \Lambda^{\mu}{}_{\nu}\,x^\nu.
  \label{eq:lec02:lorentz-transform}
\end{equation}
Invariance of \(x^2=x^\mu x_\mu\) implies
\begin{equation}
  \Lambda^{T} g\,\Lambda = g,
  \label{eq:lec02:Lorentz-condition}
\end{equation}
equivalently \(g\,\Lambda^T g = \Lambda^{-1}\).

\dfn{Timelike / lightlike / spacelike}{For any four-vector \(A^\mu\),
\begin{equation}
  A^2 \equiv A\cdot A = (A^0)^2 - \abs{\bvec{A}}^2.
  \label{eq:lec02:A-square}
\end{equation}
If \(A^2>0\), \(A^\mu\) is \emph{timelike}. If \(A^2=0\), it is
\emph{lightlike}. If \(A^2<0\), it is \emph{spacelike}.}

\begin{remark}[One \(\Lambda\) per Lorentz index]
Tensors transform with one Lorentz matrix \(\Lambda\) for each Lorentz
index. For example, a rank-2 tensor \(A^{\mu\nu}\) transforms as
\begin{equation}
  A^{\mu'\nu'} = \Lambda^{\mu}{}_{\sigma}\Lambda^{\nu}{}_{\rho}\,A^{\sigma\rho},
  \label{eq:lec02:tensor-transform}
\end{equation}
and a rank-3 tensor similarly with three factors of \(\Lambda\).
\end{remark}

%%──────────────────────────────────────────────────────────────────
\subsection*{Natural units}

\dfn{Natural units}{We often work in units where
\begin{equation}
  \hbar = c = 1.
  \label{eq:lec02:natural-units}
\end{equation}
Then
\begin{equation}
  [\text{length}] = [\text{time}] = [E^{-1}] = [m^{-1}].
  \label{eq:lec02:natural-units-dimensions}
\end{equation}}

\begin{remark}[Electron mass example]
The lecture emphasized that in natural units, a mass can be expressed as
an energy (rest energy) and also as an inverse length scale (inverse
Compton wavelength). For the electron, \(m_e\simeq 0.511\,\mathrm{MeV}\),
and this corresponds to a characteristic inverse length scale.
\end{remark}

%%──────────────────────────────────────────────────────────────────
\subsection*{To Remember}

\begin{itemize}
  \item Classical field theory treats fields \(\phi(x,t)\) as generalized
    coordinates; the action is \(S=\int dt\,dx\,\lagr\).
  \item Under the assumptions (locality, reality, at most first
    derivatives), \(\delta S=0\) gives
    \(\frac{\partial \lagr}{\partial \phi}
    = \partial_t(\partial \lagr/\partial \dot\phi)
    + \partial_x(\partial \lagr/\partial \phi')\).
  \item \(\lagr=\frac12\dot\phi^2 - v^2(\phi')^2\) gives the wave equation
    \(\ddot\phi=v^2\phi''\) and plane waves with \(\omega=kv\).
  \item Adding \(-m^2\phi^2\) yields the Klein--Gordon equation with
    \(\omega^2=(kv)^2+m^2\).
  \item Interactions \(\phi^n\) with \(n\ge 3\) generically lead to
    nonlinear equations, motivating perturbation theory.
  \item Lorentz invariance is made manifest with four-vectors, the metric
    \(g_{\mu\nu}=\mathrm{diag}(1,-1,-1,-1)\), and contractions of one
    upper with one lower index.
  \item The box operator \(\Box=\partial^\mu\partial_\mu\) packages the
    wave operator, and \((\Box+m^2)\phi=0\) is the covariant KG equation.
  \item In natural units \(\hbar=c=1\), energies, masses, inverse
    lengths, and inverse times are interchangeable as units.
\end{itemize}

%%──────────────────────────────────────────────────────────────────
\subsection*{Further Reading}

The lecturer explicitly recommended the following:

\begin{itemize}
  \item \textbf{Griffiths, \textit{Introduction to Elementary Particles}, Chapter~3}
    — Relativistic kinematics and four-vector notation.  The lecturer suggested
    this as a reference for the special-relativity review and the $g_{\mu\nu}$
    conventions used in this course.

  \item \textbf{Peskin \& Schroeder, \textit{An Introduction to QFT}, \S2.2}
    — ``The Klein--Gordon Field as Harmonic Oscillators'': covers the classical
    field Lagrangian, the Euler--Lagrange equation for fields, and the
    Klein--Gordon equation.  The lecturer described P\&S as ``one of the Bibles
    of particle physics''.

  \item \textbf{Goldstein, Poole \& Safko, \textit{Classical Mechanics} (3rd ed.),
    Chapter~13} — Lagrangian field theory from a classical-mechanics perspective,
    bridging the gap between particle mechanics (generalised coordinates $q_i(t)$)
    and field theory (the field $\phi(x,t)$ as a continuous generalised coordinate).
\end{itemize}

Additional recommended reading:

\begin{itemize}
  \item \textbf{Tong, \textit{Quantum Field Theory} (lecture notes), Chapter~1}
    (\href{https://www.damtp.cam.ac.uk/user/tong/qft.html}{\texttt{damtp.cam.ac.uk/user/tong/qft.html}},
    free): ``Classical Field Theory'' — Lagrangian densities, Noether's theorem,
    energy--momentum tensor.  Closely parallels the material of this lecture.

  \item \textbf{Halzen \& Martin, \textit{Quarks and Leptons}, Chapter~1}
    — Four-vector notation and relativistic kinematics in the particle-physics
    style, with worked examples of Lorentz invariants.
\end{itemize}
